\documentclass[reqno]{amsart}
\usepackage[T1]{fontenc}
\usepackage{amsmath,amssymb,amsthm}
\usepackage{xr-hyper}
\usepackage[hidelinks]{hyperref}
\externaldocument[sm-]{20261125T000000--slope-method__workbook}
\externaldocument[st-]{20261124T214500--supersingular-theta-divisibility__workbook}
\externaldocument[ip-]{20261124T214600--the-isotropic-plane__workbook}
\externaldocument[ch-]{20261125T003000--the-chain-and-the-square-root__workbook}
\externaldocument[hn-]{20261125T013500--hecke-nondegeneracy-on-the-tree__workbook}
\externaldocument[ho-]{20261125T016500--hecke-operators-and-theta-lifts__workbook}

% Chapter number; \ref{chap} gives it.
\newcounter{chap}
% Chapter numbers, generated from 20261125T012000--grh-workbook__guide.org
% by scripts/make-chapters.py.  Do not edit: change the order of the
% index and rerun the script (./preview-site.sh and CI do this).
%
% Usage in a sheet, after \newcounter{chap}:
%   % Chapter numbers, generated from 20261125T012000--grh-workbook__guide.org
% by scripts/make-chapters.py.  Do not edit: change the order of the
% index and rerun the script (./preview-site.sh and CI do this).
%
% Usage in a sheet, after \newcounter{chap}:
%   % Chapter numbers, generated from 20261125T012000--grh-workbook__guide.org
% by scripts/make-chapters.py.  Do not edit: change the order of the
% index and rerun the script (./preview-site.sh and CI do this).
%
% Usage in a sheet, after \newcounter{chap}:
%   \input{chapters}\chapnumber{<sheet ID>}
\makeatletter
\providecommand\chapters@set[2]{\@namedef{chapters@#1}{#2}}
\chapters@set{20261125T001000}{1}% zeros-to-ramp-sums
\chapters@set{20261125T001100}{2}% inert-primes-and-estimate-b
\chapters@set{20261125T005000}{3}% galois-fourier-and-averaging
\chapters@set{20261125T005500}{4}% columns-from-a-sphere
\chapters@set{20261125T000000}{5}% slope-method
\chapters@set{20261126T030100}{6}% a-twist-for-each-prime
\chapters@set{20261125T003000}{7}% the-chain-and-the-square-root
\chapters@set{20261125T010000}{8}% rank-lemma-and-zero-estimate
\chapters@set{20261125T010500}{9}% poincare-series-and-near-orthogonality
\chapters@set{20261124T223000}{10}% elliptic-curves-for-the-proof
\chapters@set{20261125T009500}{11}% fock-model-and-theta-intertwiners
\chapters@set{20261124T223100}{12}% complex-multiplication-and-inert-reduction
\chapters@set{20261125T008000}{13}% one-cm-determinant
\chapters@set{20261124T214500}{14}% supersingular-theta-divisibility
\chapters@set{20261124T214600}{15}% the-isotropic-plane
\chapters@set{20261125T013700}{16}% pairings-on-a-formal-group-of-height-two
\chapters@set{20261125T013500}{17}% hecke-nondegeneracy-on-the-tree
\chapters@set{20261126T030200}{18}% the-bound-is-attained
\chapters@set{20261125T009000}{19}% kronecker-limit-and-chowla-selberg
\chapters@set{20261125T002000}{20}% definite-theta-lifts
\chapters@set{20261125T015500}{21}% siegel-weil-in-the-smallest-case
\chapters@set{20261125T002100}{22}% definite-theta-lifts-and-norms
\chapters@set{20261125T016500}{23}% hecke-operators-and-theta-lifts
\chapters@set{20261125T016700}{24}% divisibility-along-the-chain
\chapters@set{20261125T004000}{25}% denominators-and-sha
\chapters@set{20261125T004100}{26}% tate-twists-and-sha
\chapters@set{20261126T020430}{27}% the-minus-part-by-gauss-sums
\chapters@set{20261125T014100}{28}% congruence-modules
\chapters@set{20261125T014500}{29}% eisenstein-congruences-and-selmer-classes
\chapters@set{20261125T014700}{30}% pseudocharacters-and-the-length-bound
\chapters@set{20261125T011000}{31}% galois-representations-black-box
\chapters@set{20261126T030300}{32}% the-data
\chapters@set{20261126T020210}{33}% bessel-poincare-series-and-their-gram-operators
\chapters@set{20261126T030400}{34}% the-seesaw-identity
\chapters@set{20261126T030500}{35}% divisibility-of-the-composite-lift
\chapters@set{20261126T030600}{36}% theorem-a-short-forms
\chapters@set{20261126T030700}{37}% the-determinant-of-a-chain
\chapters@set{20261126T030800}{38}% the-cost-of-one-lift
\chapters@set{20261126T030900}{39}% denominators-of-projections-to-lifts
\chapters@set{20261126T031000}{40}% twists-with-few-congruences
\chapters@set{20261126T031100}{41}% theorem-c-and-the-main-inequality
\providecommand\chapnumber[1]{%
  \@ifundefined{chapters@#1}%
    {\PackageError{chapters}{Sheet #1 is not listed in the index}%
      {List it in 20261125T012000--grh-workbook__guide.org\MessageBreak
       and run scripts/make-chapters.py.}}%
    {\edef\thechap{\@nameuse{chapters@#1}}}}
\makeatother
\chapnumber{<sheet ID>}
\makeatletter
\providecommand\chapters@set[2]{\@namedef{chapters@#1}{#2}}
\chapters@set{20261125T001000}{1}% zeros-to-ramp-sums
\chapters@set{20261125T001100}{2}% inert-primes-and-estimate-b
\chapters@set{20261125T005000}{3}% galois-fourier-and-averaging
\chapters@set{20261125T005500}{4}% columns-from-a-sphere
\chapters@set{20261125T000000}{5}% slope-method
\chapters@set{20261126T030100}{6}% a-twist-for-each-prime
\chapters@set{20261125T003000}{7}% the-chain-and-the-square-root
\chapters@set{20261125T010000}{8}% rank-lemma-and-zero-estimate
\chapters@set{20261125T010500}{9}% poincare-series-and-near-orthogonality
\chapters@set{20261124T223000}{10}% elliptic-curves-for-the-proof
\chapters@set{20261125T009500}{11}% fock-model-and-theta-intertwiners
\chapters@set{20261124T223100}{12}% complex-multiplication-and-inert-reduction
\chapters@set{20261125T008000}{13}% one-cm-determinant
\chapters@set{20261124T214500}{14}% supersingular-theta-divisibility
\chapters@set{20261124T214600}{15}% the-isotropic-plane
\chapters@set{20261125T013700}{16}% pairings-on-a-formal-group-of-height-two
\chapters@set{20261125T013500}{17}% hecke-nondegeneracy-on-the-tree
\chapters@set{20261126T030200}{18}% the-bound-is-attained
\chapters@set{20261125T009000}{19}% kronecker-limit-and-chowla-selberg
\chapters@set{20261125T002000}{20}% definite-theta-lifts
\chapters@set{20261125T015500}{21}% siegel-weil-in-the-smallest-case
\chapters@set{20261125T002100}{22}% definite-theta-lifts-and-norms
\chapters@set{20261125T016500}{23}% hecke-operators-and-theta-lifts
\chapters@set{20261125T016700}{24}% divisibility-along-the-chain
\chapters@set{20261125T004000}{25}% denominators-and-sha
\chapters@set{20261125T004100}{26}% tate-twists-and-sha
\chapters@set{20261126T020430}{27}% the-minus-part-by-gauss-sums
\chapters@set{20261125T014100}{28}% congruence-modules
\chapters@set{20261125T014500}{29}% eisenstein-congruences-and-selmer-classes
\chapters@set{20261125T014700}{30}% pseudocharacters-and-the-length-bound
\chapters@set{20261125T011000}{31}% galois-representations-black-box
\chapters@set{20261126T030300}{32}% the-data
\chapters@set{20261126T020210}{33}% bessel-poincare-series-and-their-gram-operators
\chapters@set{20261126T030400}{34}% the-seesaw-identity
\chapters@set{20261126T030500}{35}% divisibility-of-the-composite-lift
\chapters@set{20261126T030600}{36}% theorem-a-short-forms
\chapters@set{20261126T030700}{37}% the-determinant-of-a-chain
\chapters@set{20261126T030800}{38}% the-cost-of-one-lift
\chapters@set{20261126T030900}{39}% denominators-of-projections-to-lifts
\chapters@set{20261126T031000}{40}% twists-with-few-congruences
\chapters@set{20261126T031100}{41}% theorem-c-and-the-main-inequality
\providecommand\chapnumber[1]{%
  \@ifundefined{chapters@#1}%
    {\PackageError{chapters}{Sheet #1 is not listed in the index}%
      {List it in 20261125T012000--grh-workbook__guide.org\MessageBreak
       and run scripts/make-chapters.py.}}%
    {\edef\thechap{\@nameuse{chapters@#1}}}}
\makeatother
\chapnumber{<sheet ID>}
\makeatletter
\providecommand\chapters@set[2]{\@namedef{chapters@#1}{#2}}
\chapters@set{20261125T001000}{1}% zeros-to-ramp-sums
\chapters@set{20261125T001100}{2}% inert-primes-and-estimate-b
\chapters@set{20261125T005000}{3}% galois-fourier-and-averaging
\chapters@set{20261125T005500}{4}% columns-from-a-sphere
\chapters@set{20261125T000000}{5}% slope-method
\chapters@set{20261126T030100}{6}% a-twist-for-each-prime
\chapters@set{20261125T003000}{7}% the-chain-and-the-square-root
\chapters@set{20261125T010000}{8}% rank-lemma-and-zero-estimate
\chapters@set{20261125T010500}{9}% poincare-series-and-near-orthogonality
\chapters@set{20261124T223000}{10}% elliptic-curves-for-the-proof
\chapters@set{20261125T009500}{11}% fock-model-and-theta-intertwiners
\chapters@set{20261124T223100}{12}% complex-multiplication-and-inert-reduction
\chapters@set{20261125T008000}{13}% one-cm-determinant
\chapters@set{20261124T214500}{14}% supersingular-theta-divisibility
\chapters@set{20261124T214600}{15}% the-isotropic-plane
\chapters@set{20261125T013700}{16}% pairings-on-a-formal-group-of-height-two
\chapters@set{20261125T013500}{17}% hecke-nondegeneracy-on-the-tree
\chapters@set{20261126T030200}{18}% the-bound-is-attained
\chapters@set{20261125T009000}{19}% kronecker-limit-and-chowla-selberg
\chapters@set{20261125T002000}{20}% definite-theta-lifts
\chapters@set{20261125T015500}{21}% siegel-weil-in-the-smallest-case
\chapters@set{20261125T002100}{22}% definite-theta-lifts-and-norms
\chapters@set{20261125T016500}{23}% hecke-operators-and-theta-lifts
\chapters@set{20261125T016700}{24}% divisibility-along-the-chain
\chapters@set{20261125T004000}{25}% denominators-and-sha
\chapters@set{20261125T004100}{26}% tate-twists-and-sha
\chapters@set{20261126T020430}{27}% the-minus-part-by-gauss-sums
\chapters@set{20261125T014100}{28}% congruence-modules
\chapters@set{20261125T014500}{29}% eisenstein-congruences-and-selmer-classes
\chapters@set{20261125T014700}{30}% pseudocharacters-and-the-length-bound
\chapters@set{20261125T011000}{31}% galois-representations-black-box
\chapters@set{20261126T030300}{32}% the-data
\chapters@set{20261126T020210}{33}% bessel-poincare-series-and-their-gram-operators
\chapters@set{20261126T030400}{34}% the-seesaw-identity
\chapters@set{20261126T030500}{35}% divisibility-of-the-composite-lift
\chapters@set{20261126T030600}{36}% theorem-a-short-forms
\chapters@set{20261126T030700}{37}% the-determinant-of-a-chain
\chapters@set{20261126T030800}{38}% the-cost-of-one-lift
\chapters@set{20261126T030900}{39}% denominators-of-projections-to-lifts
\chapters@set{20261126T031000}{40}% twists-with-few-congruences
\chapters@set{20261126T031100}{41}% theorem-c-and-the-main-inequality
\providecommand\chapnumber[1]{%
  \@ifundefined{chapters@#1}%
    {\PackageError{chapters}{Sheet #1 is not listed in the index}%
      {List it in 20261125T012000--grh-workbook__guide.org\MessageBreak
       and run scripts/make-chapters.py.}}%
    {\edef\thechap{\@nameuse{chapters@#1}}}}
\makeatother
\chapnumber{20261125T016700}
\renewcommand{\thesection}{\thechap.\arabic{section}}

\theoremstyle{definition}
\newtheorem{problem}{Problem}[section]
\theoremstyle{plain}
\newtheorem{theorem}{Theorem}[section]

\newcommand{\Z}{\mathbb{Z}}
\newcommand{\Q}{\mathbb{Q}}
\newcommand{\Oc}{\mathcal{O}}
\newcommand{\Ec}{\mathcal{E}}
\newcommand{\Kc}{\mathcal{K}}
\newcommand{\Int}{\mathrm{Int}}
\newcommand{\Price}{\mathrm{Pr}}
\newcommand{\len}{\operatorname{length}}
\newcommand{\gain}{\operatorname{gain}}
\newcommand{\ord}{\operatorname{ord}}
\newcommand{\diag}{\operatorname{diag}}
\newcommand{\qbinom}[3]{\genfrac{[}{]}{0pt}{}{#1}{#2}_{#3}}

\setlength{\emergencystretch}{2em}

\title{Divisibility along the chain}
\author{}
\date{November 25, 2026}

\begin{document}
\maketitle
\refstepcounter{chap}\label{chap}

\paragraph{Prerequisites.}
Formal power series over a commutative ring.  Finite extensions of \(\Q_p\), their valuations, and lattices over their rings of integers; the length of a module of finite length.  Gaussian binomial coefficients \(\qbinom{n}{l}{Q}\) \cite[Ch.~I, \S2, Example~3]{Mac}.  From Chapter~\ref{sm-chap}, \emph{The slope method}: Theorem~\ref{sm-thm:one}.  From Chapter~\ref{ch-chap}, \emph{The chain and the square root}: Section~\ref{ch-sec:chain}.  From Chapter~\ref{hn-chap}, \emph{Hecke nondegeneracy on the tree}: the loss, defined in Section~\ref{hn-sec:tree}.  From Chapter~\ref{ho-chap}, \emph{Hecke operators and theta lifts}: Section~\ref{ho-sec:split}, in particular Problems~\ref{ho-pr:split-el} and~\ref{ho-pr:split-chain}.

Theorem~\ref{ch-thm:chain} adds gains that hold on the whole of each lattice.  In the argument a lift gains more only on a sublattice spanned by images of a Hecke operator, and Chapter~\ref{hn-chap} measures the colength of such a sublattice.  How does a gain on a sublattice at the start of a chain become a gain of the composite?  The Hecke operator has to be moved along the chain, by Problem~\ref{ho-pr:split-chain}, and the coefficients that appear are complete symmetric functions.  Their denominators are the price of the move.  We follow the reconstruction of Adriazola \cite[\S10]{Adr}, with one Hecke operator in place of several.

\section{Elementary and complete symmetric functions}\label{sec:sym}

Let \(x = (x_1, \dots, x_n)\) be a finite family in a commutative ring.  The \emph{elementary} and \emph{complete symmetric functions} \(e_l(x)\) and \(h_l(x)\) are the sums of the products \(x_{i_1} \cdots x_{i_l}\) over \(i_1 < \dots < i_l\) and over \(i_1 \le \dots \le i_l\) respectively, with \(e_0 = h_0 = 1\).  Their generating series are
\[
  E_x(X) := \prod_{i} (1 + x_i X) = \sum_{l \ge 0} e_l(x) X^l , \qquad
  H_x(X) := \prod_{i} (1 - x_i X)^{-1} = \sum_{l \ge 0} h_l(x) X^l .
\]

\begin{problem}\label{pr:EH}
\(E_x(X) H_x(-X) = 1\), that is, \(\sum_{l=0}^r (-1)^l e_{r-l}(x)\, h_l(x) = 0\) for \(r \ge 1\).  For \(x = (a, b)\), \(h_2(x) = a^2 + ab + b^2\).
\end{problem}

This is \cite[Ch.~I, (2.6)]{Mac}.  For families \(x\) and \(y\) write \(x \sqcup y\) for the concatenated family.

\begin{problem}\label{pr:union}
\[
  e_r(x \sqcup y) = \sum_{l=0}^r e_l(x)\, e_{r-l}(y) , \qquad h_r(x \sqcup y) = \sum_{l=0}^r h_l(x)\, h_{r-l}(y) .
\]
\end{problem}

The families that occur are geometric progressions.

\begin{problem}\label{pr:geom}
For \(t \ge 0\),
\[
  e_l(1, Q, \dots, Q^t) = Q^{l(l-1)/2} \qbinom{t+1}{l}{Q} , \qquad h_l(1, Q, \dots, Q^t) = \qbinom{t+l}{l}{Q} .
\]
For instance \(h_2(1, Q) = 1 + Q + Q^2\).
\end{problem}

These are \cite[Ch.~I, \S2, Example~3]{Mac}.  Let \(q\) be invertible, \(Q := q^2\), \(g \ge 1\) and \(t := g - 1\).  A \emph{block} is a family
\begin{equation}\label{eq:block}
  x = m\,(q^t, q^{t-2}, \dots, q^{-t})
\end{equation}
of \(g\) terms with ratio \(Q^{-1}\).  These are the families of Problem~\ref{ho-pr:split-el}, where \(q^t\) is written \(Q^{t/2}\).

\begin{problem}\label{pr:sym-block}
For the block~\eqref{eq:block}, \(e_l(x) = m^l q^{-l(g-l)} \qbinom{g}{l}{Q}\) and \(h_l(x) = m^l q^{-lt} \qbinom{t+l}{l}{Q}\).
\end{problem}

The first formula is Problem~\ref{ho-pr:split-el}.

\section{Valuations}\label{sec:val}

Let \(k\) be a finite extension of \(\Q_p\), with ring of integers \(\Oc\) and valuation \(v\) normalized by \(v(p) = 1\).

\begin{problem}\label{pr:gauss-unit}
\(\qbinom{n}{l}{Q}\) is a polynomial in \(Q\) with integer coefficients and constant term \(1\), for \(0 \le l \le n\).  Hence \(v\bigl(\qbinom{n}{l}{Q}\bigr) = 0\) if \(Q \in k\) and \(v(Q) > 0\).
\end{problem}

\begin{problem}\label{pr:block-val}
Let \(a \in k^\times\) and \(Q \in k\) with \(v(Q) = \nu > 0\).  Then
\begin{align*}
  v\bigl(e_l(a, aQ, \dots, aQ^t)\bigr) &= l\,v(a) + \nu\,\frac{l(l-1)}{2} && (0 \le l \le t+1) , \\
  v\bigl(h_l(a, aQ, \dots, aQ^t)\bigr) &= l\,v(a) && (l \ge 0) .
\end{align*}
\end{problem}

So the elementary functions of a block become more divisible with \(l\), and the complete functions do not.

\begin{problem}\label{pr:sym-val}
Let \(x\) be the block~\eqref{eq:block} with \(m, q \in k^\times\) and \(v(q) = \nu > 0\).  Then \(v(h_l(x)) = l\,(v(m) - t\nu)\) and \(v(e_l(x)) = l\,(v(m) - t\nu) + \nu\, l(l-1)\) for \(0 \le l \le g\).  In particular every \(h_l(x)\) lies in \(\Oc\) if and only if \(v(m) \ge t\nu\).
\end{problem}

For a union of blocks only a bound survives.

\begin{problem}\label{pr:union-val}
Let \(Z\) be a finite family in \(k\) whose least valuation is \(\mu\).  Then \(v(h_l(Z)) \ge l\mu\), with equality if exactly one member of \(Z\) has valuation \(\mu\).  For \(Z = (1, 1)\) and \(Z = (1, -1)\) one has \(h_l(Z) = l + 1\) and \(h_l(Z) = (1 + (-1)^l)/2\).
\end{problem}

\begin{problem}\label{pr:two-blocks}
Let \(Z = x_1 \sqcup x_2\), where \(x_i = m_i\,(q, q^{-1})\) are blocks with \(g = 2\), and let \(v(q) = \nu > 0\) and \(v(m_1) = v(m_2) = 0\).  Then \(h_1(Z) = (m_1 + m_2)(q + q^{-1})\), so \(v(h_1(Z)) = v(m_1 + m_2) - \nu\), which exceeds \(-\nu\) exactly when \(m_1 \equiv -m_2\) modulo the maximal ideal of \(\Oc\).
\end{problem}

\section{Moving a Hecke operator along a chain}\label{sec:move}

We use the setting of Problem~\ref{ho-pr:split-chain}.  Let \(F_0, \dots, F_m\) be finite-dimensional \(k\)-vector spaces, let \(K_j \colon F_j \to F_{j+1}\) be linear maps, and let \(\Kc_j := K_{j-1} \cdots K_0\), with \(\Kc_0 := 1\).  Fix \(r \ge 1\).  For each \(j\) let \(\Ec^{(j)}_0 = 1, \Ec^{(j)}_1, \dots, \Ec^{(j)}_r\) be operators on \(F_j\), and let \(x_1, \dots, x_m\) be finite families in \(k\) with
\begin{equation}\label{eq:R}
  \Ec^{(j+1)}_i\, K_j = K_j \sum_{l=0}^i e_l(x_{j+1})\, \Ec^{(j)}_{i-l} \qquad (0 \le i \le r,\ 0 \le j < m) .
\end{equation}
For lifts with split kernels, Section~\ref{ho-sec:split} gives~\eqref{eq:R} with \(x_{j+1}\) a block~\eqref{eq:block}, up to a unit that we take to be \(1\) \cite[Props.~10.8 and~10.13]{Adr}.  Let \(Z_j := x_1 \sqcup \dots \sqcup x_j\).  In~\eqref{eq:R} the operators at the target are expressed through those at the source.  We need the converse.

\begin{problem}\label{pr:r-one}
Let \(r = 1\).  Then \(K_0 \Ec^{(0)}_1 = \bigl(\Ec^{(1)}_1 - e_1(x_1)\bigr) K_0\), and \(\Kc_2 \Ec^{(0)}_1 = \bigl(\Ec^{(2)}_1 - h_1(Z_2)\bigr) \Kc_2\).
\end{problem}

\begin{problem}\label{pr:r-two}
Let \(r = 2\).  Then \(K_0 \Ec^{(0)}_2 = \bigl(\Ec^{(1)}_2 - h_1(x_1)\, \Ec^{(1)}_1 + h_2(x_1)\bigr) K_0\).
\end{problem}

\begin{problem}\label{pr:forward}
\(\Ec^{(j)}_i\, \Kc_j = \Kc_j \sum_{l=0}^i e_l(Z_j)\, \Ec^{(0)}_{i-l}\) for \(0 \le i \le r\).
\end{problem}

Problems~\ref{pr:forward} and~\ref{pr:EH} give the following, which is Problem~\ref{ho-pr:split-chain}.

\begin{theorem}\label{thm:move}
For \(1 \le j \le m\),
\(\displaystyle \Kc_j\, \Ec^{(0)}_r = \sum_{l=0}^r (-1)^l h_l(Z_j)\, \Ec^{(j)}_{r-l}\, \Kc_j\).
\end{theorem}

The relation for one lift has the elementary functions of a block, whose valuations grow with \(l\) (Problem~\ref{pr:sym-val}).  The inverse relation has the complete functions, which may have denominators.  The operators \(\Ec^{(j)}_i\) are normalized forms of the sums \(T^{(j)}_i\) over edges of type \(i\) of Section~\ref{ho-sec:split}: when \(F_j\) consists of functions of lattices in a space of dimension \(N_j\), with \(N_0 < N_1 < \dots < N_m\), one takes \(\Ec^{(j)}_i = q^{-i(N_j - i)}\, T^{(j)}_i\) \cite[\S10.3]{Adr}.

\begin{problem}\label{pr:raw}
With this normalization, Theorem~\ref{thm:move} reads
\[
  \sum_{l=0}^r (-1)^l h_l(Z_j)\, q^{l(N_j - 2r + l)}\, T^{(j)}_{r-l}\, \Kc_j = q^{r(N_j - N_0)}\, \Kc_j\, T^{(0)}_r .
\]
\end{problem}

\begin{problem}\label{pr:integral}
Suppose that \(Z_j\) is a union of blocks~\eqref{eq:block} with the same \(t\), that \(v(q) > 0\), that \(v(m) \ge 0\) for each block, and that \(N_j \ge 2r + t - 1\).  Then every coefficient \(h_l(Z_j)\, q^{l(N_j - 2r + l)}\) in Problem~\ref{pr:raw} lies in \(\Oc\).
\end{problem}

\section{A gain on a sublattice}\label{sec:gain}

For each \(j\) let \(\Int_j \subset F_j\) be an \(\Oc\)-lattice of full rank.  For \(w \in F_j\) let \(\ord(w)\) be the least valuation of a coordinate of \(w\) in a basis of \(\Int_j\), with \(\ord(0) = \infty\).  For a linear map \(A \colon F_i \to F_j\) let \(\gain(A)\) be the least value of \(\ord(Ab)\) over the members \(b\) of a basis of \(\Int_i\).  For \(k = \Q_p\), \(\gain(A)\) is the largest integer \(g\) with \(A(\Int_i) \subseteq p^g \Int_j\); so \(\gain(A) \log p\) plays at \(p\) the role of the divisibility \(-h_{\mathrm{fin}}\) of Section~\ref{ch-sec:chain}.

\begin{problem}\label{pr:gain}
\(\ord\) and \(\gain\) do not depend on the bases.  \(\ord(Aw) \ge \ord(w) + \gain(A)\), \(\gain(AB) \ge \gain(A) + \gain(B)\), \(\gain(A + B) \ge \min(\gain(A), \gain(B))\), and \(\gain(cA) = v(c) + \gain(A)\) for \(c \in k^\times\).
\end{problem}

Assume~\eqref{eq:R}.  For \(1 \le j \le m\) and \(0 \le i \le r\) let
\[
  \gamma_j(i) := \gain\bigl(\Ec^{(j)}_i K_{j-1}\bigr) ,
\]
so that \(\gamma_j(0) = \gain(K_{j-1})\) is the gain of the \(j\)th step on the whole lattice, and \(\gamma_j(r)\) its gain after the Hecke operator \(\Ec^{(j)}_r\) at the target.  Let \(c \in k^\times\) be such that \(Y := c\, \Ec^{(0)}_r\) maps \(\Int_0\) into itself, and let \(A := v(c)\).  The gain of the \(j\)th step on the image of \(Y\) is not \(A + \gamma_j(r)\).  By Theorem~\ref{thm:move} it is controlled by
\[
  \delta_j := \min_{0 \le l \le r} \bigl( v(h_l(Z_j)) + \gamma_j(r - l) \bigr) ,
  \qquad \Price_j := \gamma_j(r) - \delta_j .
\]
When \(\gamma_j(r) < \infty\), we call \(\Price_j\) the \emph{price} of the \(j\)th step; the term \(l = 0\) shows that \(\Price_j \ge 0\), and \(\delta_j = \gamma_j(r) - \Price_j\).

\begin{problem}\label{pr:one-step}
For \(w \in F_0\) and \(1 \le j \le m\), \(\ord\bigl(\Kc_j\, \Ec^{(0)}_r w\bigr) \ge \ord\bigl(\Kc_{j-1} w\bigr) + \delta_j\).
\end{problem}

\begin{problem}\label{pr:induction}
Let \(S \subseteq \{1, \dots, m\}\), and let \(n_j\) be the number of elements of \(S\) that are at most \(j\).  For \(G \in \Int_0\) and \(0 \le j \le m\),
\[
  \ord\bigl(\Kc_j\, Y^{n_j} G\bigr) \ge \sum_{i \le j,\ i \in S} \bigl(A + \delta_i\bigr) + \sum_{i \le j,\ i \notin S} \gamma_i(0) .
\]
\end{problem}

Take \(S\) to be the set of \(j\) with \(A + \delta_j > \gamma_j(0)\), and \(G := Y^{m - n_m} F\), which lies in \(\Int_0\) because \(Y\) preserves \(\Int_0\).  Problem~\ref{pr:induction} with \(j = m\) gives the following.

\begin{theorem}\label{thm:chain-gain}
For every \(F \in \Int_0\),
\[
  \ord\bigl(\Kc_m\, Y^m F\bigr) \ge \sum_{j=1}^m \max\bigl(\gamma_j(0),\ A + \delta_j\bigr) .
\]
\end{theorem}

This is the case of one orbit of \cite[Thm.~10.7]{Adr}, in which the estimates at several primes are combined.  A step uses a factor of \(Y\) only when the factor pays for its price; the other factors are applied at the start, where they cost nothing.  The gain holds on the sublattice \(Y^m \Int_0\).  In Theorem~\ref{sm-thm:one} the gain on a sublattice is charged its colength, divided by the rank.

\begin{problem}\label{pr:colength}
If \(Y\) is injective, then the length of \(\Int_0/Y^m\Int_0\) is \(m\) times the length of \(\Int_0/Y\,\Int_0\).
\end{problem}

When \(\Int_0\) consists of integral forms and \(Y\) is the normalized Hecke operator of Chapter~\ref{hn-chap}, the values of \(Y\) at one lattice \(L\) span the module \(S(L)\) there, and the loss of Chapter~\ref{hn-chap} is its colength in \(H(L)\).

\section{Stepwise maxima overcount}\label{sec:overcount}

Theorem~\ref{thm:chain-gain}, like \cite[Thm.~10.7]{Adr}, charges each step a price; in \cite[(10.2)]{Adr} an upper bound \(\mathrm{PR}(v, j)\) for it is subtracted from the gain of the \(j\)th lift.  Can the price be ignored, so that each step gains \(\max(\gamma_j(0), A + \gamma_j(r))\) on \(Y^m\Int_0\)?  We compare two chains with the same numbers \(\gamma_j(i)\).  Let \(p\) be any prime, \(k = \Q_p\), \(m = 2\), \(r = 1\), \(F_j = \Q_p^2\) and \(\Int_j = \Z_p^2\).  Let \(K_0 = K_1 = K\) be a diagonal matrix, let \(x_1 = (a_1)\) and \(x_2 = (a_2)\) be families of one element, and let
\[
  Y := \Ec^{(0)}_1 := \diag(p, 1) , \qquad \Ec^{(j)}_1 := Y + h_1(Z_j) \quad (j = 1, 2) ,
\]
so that \(A = 0\).

\begin{problem}\label{pr:diag}
These data satisfy~\eqref{eq:R}, and \(\len \Z_p^2/Y\Z_p^2 = 1\).
\end{problem}

\begin{problem}\label{pr:good}
Let \(K = \diag(1, p)\) and \(a_1 = a_2 = p\).  Then \(\gamma_j(0) = 0\), \(\gamma_j(1) = 1\) and \(\Price_j = 0\) for \(j = 1, 2\), and \(\gain(\Kc_2\, Y^2) = 2\).
\end{problem}

\begin{problem}\label{pr:bad}
Let \(K = \diag(p, 1)\), \(a_1 = p - 1\) and \(a_2 = -2p\).  Then \(\gamma_j(0) = 0\), \(\gamma_j(1) = 1\) and \(\Price_j = 1\) for \(j = 1, 2\), and \(\gain(\Kc_2\, Y^2) = 0\).
\end{problem}

In both chains each step, followed by the Hecke operator at its target, gains \(1\), and the sum of the stepwise maxima is \(2\).  In the second chain the composite gains nothing on \(Y^2\Int_0\): the Hecke operator at the target of each step is not the image of the one at the source, and \(h_1(Z_j)\) is a unit.  Theorem~\ref{thm:chain-gain} gives the exact gain in both cases.

The theorem leaves two inputs open: the gains \(\gamma_j(i)\) of a single lift followed by a Hecke operator at its target, which come from the divisibility of theta kernels (Chapters~\ref{st-chap} and~\ref{ip-chap}), and the colength of \(Y\,\Int_0\), which is the subject of Chapter~\ref{hn-chap}.  The price couples them, through the least valuation in the union \(Z_j\) of the blocks met so far (Problems~\ref{pr:sym-val} and~\ref{pr:union-val}).

\subsection*{Exercises}

\begin{enumerate}
\item The Newton polygon of \(\prod_{i=0}^t (1 + aQ^iX)\) over \(k\), with \(v(Q) > 0\), has the vertices \(\bigl(l, l\,v(a) + v(Q)\,l(l-1)/2\bigr)\) for \(0 \le l \le t+1\).
\item Suppose that~\eqref{eq:R} holds with \(u_{j+1}^i K_j\) in place of \(K_j\) on the right, for units \(u_j \in \Oc^\times\).  In terms of \(E_j(X) := \sum_{i \le r} \Ec^{(j)}_i X^i\) this reads \(E_{j+1}(X) K_j = K_j\, E_{x_{j+1}}(u_{j+1}X)\, E_j(u_{j+1}X)\) modulo \(X^{r+1}\); it is the relation of Section~\ref{ho-sec:split} with the family \(u_{j+1}x_{j+1}\), whose product \(\prod_a (1 + aX)\) is the \(P(X)\) written there.  Then Theorem~\ref{thm:move} holds with \(\Kc_j\Ec^{(0)}_r\) multiplied by \((u_1 \cdots u_j)^r\), and with \(Z_j\) replaced by the family \(Z'_j\) of the \(u_l u_{l+1} \cdots u_j\, a\) for \(l \le j\) and \(a\) in \(x_l\) \cite[Prop.~10.8]{Adr}.  Theorem~\ref{thm:chain-gain} holds with \(\delta_j\) computed from \(Z'_j\).  These agree with the old \(\delta_j\) when \(Z_j\) has exactly one member of least valuation (Problem~\ref{pr:union-val}), but not in general: for \(p\) odd, \(x_1 = (1)\), \(x_2 = (-1)\) and \(u_1 = -1\), one has \(h_1(Z_2) = 0\) while \(h_1(Z'_2)\) is a unit.
\end{enumerate}

\section*{Hints}

\subsection*{\S\ref{sec:sym}}
\begin{description}
\item[\ref{pr:EH}] Compare the product formulas.
\item[\ref{pr:union}] \(E_{x \sqcup y} = E_x E_y\).
\item[\ref{pr:geom}] Induction on \(t\), with the rule \(\qbinom{n}{l}{Q} = \qbinom{n-1}{l-1}{Q} + Q^l \qbinom{n-1}{l}{Q}\) and its companion \(\qbinom{n}{l}{Q} = Q^{n-l}\qbinom{n-1}{l-1}{Q} + \qbinom{n-1}{l}{Q}\).
\item[\ref{pr:sym-block}] Problem~\ref{pr:geom} for the family \(mq^{-t}(1, Q, \dots, Q^t)\); \(e_l\) and \(h_l\) are homogeneous of degree \(l\).
\end{description}

\subsection*{\S\ref{sec:val}}
\begin{description}
\item[\ref{pr:gauss-unit}] The first rule in the hint to Problem~\ref{pr:geom}.
\item[\ref{pr:block-val}] Problems~\ref{pr:geom} and~\ref{pr:gauss-unit}.
\item[\ref{pr:sym-val}] Problems~\ref{pr:sym-block} and~\ref{pr:gauss-unit}.
\item[\ref{pr:union-val}] Which monomial of degree \(l\) has the least valuation?
\item[\ref{pr:two-blocks}] Problem~\ref{pr:union}.
\end{description}

\subsection*{\S\ref{sec:move}}
\begin{description}
\item[\ref{pr:r-one}] Apply~\eqref{eq:R} with \(i = 1\) once for each step.
\item[\ref{pr:r-two}] Use~\eqref{eq:R} with \(i = 1\) to remove \(K_0\Ec^{(0)}_1\), then Problem~\ref{pr:EH}.
\item[\ref{pr:forward}] Write~\eqref{eq:R} as an identity of polynomials in \(X\) modulo \(X^{r+1}\), as in the hint to Problem~\ref{ho-pr:split-chain}, and use Problem~\ref{pr:union}.
\item[\ref{pr:raw}] Multiply Theorem~\ref{thm:move} by \(q^{r(N_j - r)}\).
\item[\ref{pr:integral}] Problems~\ref{pr:sym-val} and~\ref{pr:union-val}; look at \(l = 1\).
\end{description}

\subsection*{\S\ref{sec:gain}}
\begin{description}
\item[\ref{pr:gain}] Write \(w\) in coordinates.
\item[\ref{pr:one-step}] Apply Theorem~\ref{thm:move} to \(w\), write \(\Kc_j = K_{j-1}\Kc_{j-1}\), and use Problem~\ref{pr:gain}.
\item[\ref{pr:induction}] Induction on \(j\).  For \(j \in S\) write \(Y^{n_j} = c\,\Ec^{(0)}_r\, Y^{n_{j-1}}\) and use Problem~\ref{pr:one-step}.
\item[\ref{pr:colength}] \(Y^i\Int_0 / Y^{i+1}\Int_0 \cong \Int_0/Y\,\Int_0\).
\end{description}

\subsection*{\S\ref{sec:overcount}}
\begin{description}
\item[\ref{pr:diag}] Diagonal matrices commute.
\item[\ref{pr:good}] \(\Kc_2\, Y^2 = p^2\).
\item[\ref{pr:bad}] Compute \(\Ec^{(1)}_1\) and \(\Ec^{(2)}_1\) first.
\end{description}

{\small\begin{thebibliography}{Mac}
\bibitem[Adr]{Adr} J.~Adriazola, \emph{The Alp\"oge proof of the Riemann hypothesis}, manuscript, revised 24 November 2026, \url{https://github.com/jimjamzola/RH-Alpoge-Reconstruction}.
\bibitem[Mac]{Mac} I.~G.~Macdonald, \emph{Symmetric Functions and Hall Polynomials}, 2nd ed., Oxford Math.\ Monogr., Clarendon Press, Oxford, 1995.
\end{thebibliography}}

\end{document}
